\documentclass[11pt,a4paper]{article}
\usepackage{xeCJK}
\usepackage{fontspec}
\usepackage{xcolor}
\usepackage{geometry}
\usepackage[protrusion=true]{microtype}
\geometry{margin=1.8cm,top=1.9cm}
\linespread{1.25}

\setmainfont{Baskerville}
\setCJKmainfont{Songti SC}

\definecolor{tipblue}{RGB}{0,95,160}
\definecolor{rulegray}{RGB}{190,190,190}

\setlength{\parindent}{0pt}
\setlength{\parskip}{0.35em}
\newcommand{\num}[1]{\makebox[1.9em][r]{#1.}\hspace{0.3em}}
\newcommand{\wline}{\noindent\rule{\linewidth}{0.4pt}}
% 中文题干：宋体粗体；提示词：蓝色小字；英文句子：Baskerville 正体
\newcommand{\prompt}[1]{\emph{\textcolor{tipblue}{\small（#1）}}}

\begin{document}
\pagestyle{plain}
\begin{center}
{\LARGE\bfseries 上海高考高分翻译背诵 \quad 64}\\[0.25em]
{\large 上海高考翻译背诵-62-25奉贤一模　2026-08-24 \quad\ 中英对照（背诵 · 答案）}\\[0.35em]
{\footnotesize 来源：微信公众号「发现之旅 速来记单词」· 2026-08-24 · 赵文瀚}
\end{center}

\vspace{0.4em}
\noindent\begin{tabular}{@{}p{0.33\textwidth}@{}p{0.34\textwidth}@{}p{0.33\textwidth}@{}}
姓名：\underline{\hspace{2.0cm}} & 日期：\underline{\hspace{2.0cm}} & 得分：\underline{\hspace{1.4cm}}\,/\enspace 4 \\
\end{tabular}

\vspace{0.5em}
\noindent{\small 中 $\to$ 英默写卷的参考答案；亦可直接用于背诵与回译自查。}

\vspace{0.8em}
\num{1}\textbf{农忙季节，这个偏远的小村子显得空荡荡的。}\prompt{desert}

\vspace{0.15em}

\noindent\hangindent=2.2em\hangafter=1\hspace{2.2em}During the busy farming season, the remote small village lies deserted.

\num{2}\textbf{上海博物馆的埃及展太火爆了，根本抢不到票。}\prompt{impossible}

\vspace{0.15em}

\noindent\hangindent=2.2em\hangafter=1\hspace{2.2em}The Egypt exhibition at the Shanghai Museum is so popular that it is virtually impossible to secure a ticket.

\num{3}\textbf{虽然这部经典老片在1998年首映，但它本月重新上映时，依然在影迷中引起了极大的反响。}\prompt{release}

\vspace{0.15em}

\noindent\hangindent=2.2em\hangafter=1\hspace{2.2em}Although the classic film was first released back in 1998, it caused a sensation among film fans upon its re-release this month.

\num{4}\textbf{这位原常深耕于科研领域的女学者，凭借自己坚持不懈的努力，历经严格的训练，最终成为了一名宇航员。}\prompt{engage}

\vspace{0.15em}

\noindent\hangindent=2.2em\hangafter=1\hspace{2.2em}Originally engaged in scientific research, the female scholar eventually became an astronaut through years of persistent effort and rigorous training.

\end{document}
